\documentclass[11pt,letterpaper]{article}
\usepackage[utf8]{inputenc}
\usepackage[T1]{fontenc}
\usepackage{lmodern}
\usepackage[margin=0.28in]{geometry}
\usepackage{hyperref}
\hypersetup{colorlinks=false, urlcolor=black, pdfborder=0 0 0}

\setlength{\parindent}{0pt}
\setlength{\parskip}{0pt}
\renewcommand{\baselinestretch}{0.82}
\setlength{\lineskip}{0pt}
\setlength{\lineskiplimit}{0pt}

\makeatletter
\renewcommand{\@listi}{%
  \leftmargin=\leftmargini
  \rightmargin=0pt
  \labelsep=4pt
  \labelwidth=16pt
  \itemindent=0pt
  \listparindent=0pt
  \topsep=0pt
  \partopsep=0pt
  \parsep=0pt
  \itemsep=-11pt
}
\let\saveditemize\itemize
\let\savedenditemize\enditemize
\renewcommand{\itemize}{\bgroup\setlength{\baselineskip}{0.65\baselineskip}\saveditemize}
\renewcommand{\enditemize}{\savedenditemize\egroup}
\makeatother

\begin{document}

{\Large \textbf{HSU, Ke-Cheng (Clio)}}\\[1pt]
(852) 6840 9753 \quad $|$ \quad \href{mailto:khsu@connect.ust.hk}{khsu@connect.ust.hk} \quad $|$ \quad \href{https://github.com/cliohsu1997}{github.com/cliohsu1997} \quad $|$ \quad \href{https://cliohsu1997.github.io}{cliohsu1997.github.io}

\vspace{16pt}

{\large \textbf{Summary}}\\[-0.5pt]
PhD candidate in Economics at HKUST, graduating June 2027. Graduated first in class (CGA 3.9/4.3; Academic Achievement Medal, top 1\%). Uses data to answer business questions and explains results clearly. Leadership and communication experience as Hall Tutor and Teaching Assistant Coordinator, with student teaching ratings of 4.5--4.7 out of 5. Hong Kong permanent resident; native Mandarin, fluent English, communicative Cantonese. Applying for the 2027 Management Trainee Programme, with Customer and Product Development as the preferred function.

\vspace{16pt}

{\large \textbf{Education}}\\[-0.5pt]
\textbf{The Hong Kong University of Science and Technology}, Hong Kong\\
\textit{Doctor of Philosophy in Economics} \hfill Expected Graduation: June 2027\\
\textit{Master of Philosophy in Economics} \hfill 2020--2022 \quad CGA: 4.1/4.3\\
\textit{BSc in Mathematics and Economics, Minor in Actuarial Mathematics} \hfill 2015--2019 \quad CGA: 3.9/4.3\\
Graduated first in class. Awards: Dean's List (2016--2019); University Scholarship (730K HKD); Academic Achievement Medal (top 1\%); Postgraduate studentship (1.5M HKD); Redbird scholarship (40K HKD).

\vspace{16pt}

{\large \textbf{Experience}}\\[-0.5pt]
\textbf{Teaching Assistant Coordinator}, Department of Economics, HKUST \hfill 2025--2026\\
\vspace{-6pt}
\begin{itemize}
  \item Evaluated teaching assistants' practice teaching, gave feedback on what to change, and helped them improve.
\end{itemize}
\textbf{Hall Tutor}, UG Hall I, HKUST \hfill 2024--present\\
\vspace{-6pt}
\begin{itemize}
  \item Helped residents work through conflicts and supported student life in a residential community.
\end{itemize}
\textbf{Lecturer \& Teaching Assistant}, Department of Economics, HKUST \hfill 2020--present\\
\vspace{-6pt}
\begin{itemize}
  \item Lecturer for MSc Math Camp (student rating 4.53/5); TA for PhD and undergraduate economics (average 4.72/5 across five courses).
\end{itemize}

\vspace{16pt}

{\large \textbf{Projects}}\\[-0.5pt]
\textbf{Return Policy Design in E-Commerce} \textit{(Job Market Paper)}\\
\vspace{-6pt}
\begin{itemize}
  \item Built and estimated a structural model of purchase, return-shipping-insurance take-up, and return decisions under monopoly pricing, using within-retailer switches between automatic provision and selective adoption of return-shipping insurance on Taobao (2017--2023; $\sim$62M RMB revenue; six apparel categories; $\sim$20--30\% return rates); scraped and OCR-digitized product attributes in Python.
  \item Developed a reproducible, version-controlled estimation pipeline on AWS and HPC/SLURM (Python, R, C++); recovered category-specific demand, consumer fit uncertainty, and return responses via GMM; accelerated Monte Carlo and optimization in C++; validated policy rankings (seller profit, consumer surplus, social welfare) with bootstrap inference.
  \item Found that a lenient policy (seller pays return shipping) suits high-margin categories, a strict policy (buyer pays) suits thin-margin ones (improving welfare by $\sim$50 pp in the thinnest), and a proposed menu (lower price with partial restocking fee vs. full price with full refund) beats both, reaching $\sim$75\% of the best attainable outcome in most categories; back-of-envelope calculation suggests $\sim$5M RMB profit improvement. The same method answers what a product team asks before a launch: how much a new product would sell, what it would contribute to profit, and whether it would cannibalize existing products.
\end{itemize}
\vspace{-3pt}
\textbf{Inventory and Pricing under Resale Price Maintenance (RPM)} \textit{(with K.\ Kawaguchi, HKUST; K.\ Kawai, U.\ Tokyo)}\\
\vspace{-6pt}
\begin{itemize}
  \item Built a dynamic oligopoly model of pricing and inventory under demand uncertainty and RPM, with equilibrium stockouts and unsold returns (35--40\% of Japanese books printed are discarded); constructed title-by-market-by-month panels from monthly bookstore POS data (tens of billions of records across 1{,}000+ stores), linked via standardized ISBNs and merged with local demographics and detailed title-level information scraped from Amazon in Python.
  \item Estimated demand and supply with instrumental variables (2SLS) and Hausman-style price instruments; validated the estimator with Monte Carlo on simulated equilibria; selected controls with LASSO and panel fixed effects; implemented iterative equilibrium simulation in C++.
  \item Found that the static model overestimates marginal cost by $\sim$20\% when ignoring dynamic demand uncertainty and inventory; uncovered stockout-driven intertitle and intertemporal substitution (not interregional) relevant for inventory and assortment decisions.
\end{itemize}

\vspace{16pt}

{\large \textbf{Skills}}\\[-0.5pt]
\textbf{Languages:} Mandarin (Native), English (Fluent), Cantonese (Communicative).\\
\textbf{Technical:} Python, R, C++, Stata, \LaTeX, AWS, HPC/SLURM; structural estimation and counterfactual simulation.\\
\textbf{References:} Kei Kawai (UTokyo; Compass Lexecon), Kohei Kawaguchi (HKUST; UTokyo Econ Consulting), Pak Hung Au (HKUST).

\end{document}
